% !TeX program = xelatex
\documentclass[11pt,a4paper]{ctexart}

\usepackage[left=2.4cm,right=2.4cm,top=2.1cm,bottom=2.1cm]{geometry}
\usepackage{amsmath,amssymb,mathtools,braket,bm}
\usepackage{xcolor}
% inline 选项用于启用 enumerate* 行内列表环境。
\usepackage[inline]{enumitem}
\usepackage{fancyhdr}

% ============================================================
% 基本信息：只需修改下面几项
% ============================================================
\newcommand{\coursename}{量子场论}
\newcommand{\assignmentnumber}{第二章习题}
\newcommand{\teachername}{\underline{乔从丰、刘泽茏}}
\newcommand{\studentname}{\underline{林照翔}}
\newcommand{\studentid}{\underline{202628001207046}}

% ============================================================
% 数学公式书写规范
% ============================================================
% 1. 行内公式使用一对美元符号。
% 2. 独立成行的公式统一使用 equation 环境，不使用反斜杠方括号形式。
% 3. 虚数单位 i、自然常数 e 和微分符号 d 均使用直立体 \mathrm。
%

% 独立公式的推荐写法：
% \begin{equation}
%   ...
% \end{equation}
%

% ============================================================
% 页面样式
% ============================================================
\setlength{\parindent}{2em}
\setlength{\parskip}{0.25em}
\setlength{\headheight}{14pt}
\pagestyle{fancy}
\fancyhf{}
\lhead{\coursename}
\rhead{\assignmentnumber}
\cfoot{\thepage}
\renewcommand{\headrulewidth}{0.4pt}

% ============================================================
% 列表参数
% ============================================================
% enumerate 是每个条目 另起一行 的普通编号列表。
% label 决定编号格式：\arabic* 依次产生 1、2、3，外加括号后显示为 (1)、(2)、(3)。
\setlist[enumerate,1]{
  label=(\arabic*),
  leftmargin=3.2em,
  itemsep=0.55em,
  topsep=0.45em
}

% inlineenum 是基于 enumerate* 定义的行内编号列表。
% 所有格式均在此处统一设置，正文中不必为每一道题重复填写参数。
% label 决定编号格式；labelsep 决定编号与本项正文之间的距离。
% itemjoin 决定相邻条目之间插入的内容。
% itemjoin* 决定最后两个条目之间插入的内容，通常与 itemjoin 保持一致。
% 内容过长时，行内列表仍会自动换行。
\newlist{inlineenum}{enumerate*}{1}
\setlist[inlineenum,1]{
  label=(\arabic*),
  labelsep=0.5em,
  itemjoin={；\quad},
  itemjoin*={；\quad}
}

% 其他常用 label 写法：
%   label=\arabic*.     -> 1. 2. 3.
%   label=(\arabic*)   -> (1) (2) (3)
%   label=\alph*)      -> a) b) c)
%   label=(\roman*)    -> (i) (ii) (iii)
%   label=\Alph*.      -> A. B. C.

% ============================================================
% 题目与回答环境
% ============================================================
% problem 的可选参数填写新版教材练习编号。
% \refstepcounter 使题目可以配合 \label 和 \ref 进行交叉引用。
% “题目一”等题目标识使用白底黑字描边框，与黑底白字的“回答”标签相区分。
\newcounter{problem}
\newenvironment{problem}[1][]
  {\refstepcounter{problem}%
   \par\medskip\noindent
   \begingroup
   \setlength{\fboxsep}{3pt}%
   \setlength{\fboxrule}{0.8pt}%
   \fbox{\bfseries 练习 #1}%
   \endgroup
   \quad\ignorespaces}
  {\par\medskip}

% answer 环境用于书写答案。
% “回答”使用黑底白字标签，保证黑白打印时仍有醒目的视觉区分；
% 答案结束处会自动添加一条黑色细分隔线。
\newenvironment{answer}
  {\par\smallskip\noindent
   \begingroup
   \setlength{\fboxsep}{4pt}%
   \colorbox{black}{\textcolor{white}{\bfseries 回答}}%
   \endgroup
   \quad\ignorespaces}
  {\par\smallskip\noindent\rule{\linewidth}{0.4pt}\medskip}

\begin{document}

\begin{center}
  {\LARGE\bfseries \coursename\quad \assignmentnumber}\par
  \vspace{0.7em}
  \small
  教师：\teachername\qquad
  姓名：\studentname\qquad
  学号：\studentid
\end{center}
\vspace{0.4em}

\begin{problem}[2.1.1]\label{prob:ch2-1}
{\small\textnormal{（第 2.1 节，第 22 页）}}\par
设拉格朗日密度为 $\mathcal L(\phi_a,\partial_\mu\phi_a)$，由最小作用量原理推导场的运动方程（教材式（2.16））：
\begin{equation}
\frac{\partial\mathcal L}{\partial\phi_a}-\partial_\mu\left(\frac{\partial\mathcal L}{\partial(\partial_\mu\phi_a)}\right)=0.
\end{equation}
\end{problem}
\begin{answer}
  作用量为
  \begin{equation}
    S = \int \mathrm{d}^4x L(\phi_a,\partial_\mu\phi_a),
  \end{equation}
  作用量的变分为
  \begin{equation}
    \begin{aligned}
      \delta S
      &=\int \mathrm{d}^4x\left(\frac{\partial \mathcal{L} }{\partial \phi_a } \delta\phi_a + \frac{\partial \mathcal{L} }{\partial \left(\partial_\mu \phi_a \right) }  \partial_\mu\left(\delta\phi_a \right) \right) \\
      &=\int \mathrm{d}^4x \left(\frac{\partial \mathcal{L} }{\partial \phi_a } - \partial_\mu \frac{\partial \mathcal{L} }{\partial\left(\partial_\mu\phi_a \right) }  \right) \delta \phi_a + \int \mathrm{d}^4x \partial_\mu \left(\frac{\partial \mathcal{L} }{\partial\left(\partial_\mu\phi_a \right) } \delta \phi_a \right),
    \end{aligned}
  \end{equation}
  上式中散度项的四维体积分可化为边界上的三维积分，应为零。由最小作用量原理，$\delta S=0 $，再由 $\delta\phi_a $ 的任意性，得到：
  \begin{equation}
    \frac{\partial \mathcal{L} }{\partial \phi_a } - \partial_\mu \frac{\partial \mathcal{L} }{\partial\left(\partial_\mu\phi_a \right) }
    =0.
  \end{equation} 
\end{answer}
\clearpage
\begin{problem}[2.2.1]\label{prob:ch2-2}
{\small\textnormal{（第 2.2 节，第 28 页）}}\par
查找文献阅读关于庞加莱（Poincar\'e）变换与相应的守恒量部分。
\end{problem}
\begin{answer}
% 在此填写答案；当前保持空白。

\end{answer}
\clearpage
\begin{problem}[2.3.1]\label{prob:ch2-3}
{\small\textnormal{（第 2.3 节，第 30 页）}}\par
推导场的哈密顿正则运动方程（教材式（2.67）中的两个方程）：
\begin{equation}
\begin{aligned}
\dot\phi_a&=\frac{\partial\mathcal H}{\partial\pi_a}
-\nabla\cdot\frac{\partial\mathcal H}{\partial(\nabla\pi_a)},\\
\dot\pi_a&=-\frac{\partial\mathcal H}{\partial\phi_a}
+\nabla\cdot\frac{\partial\mathcal H}{\partial(\nabla\phi_a)}.
\end{aligned}
\end{equation}
其中 $\phi_a$ 为场分量，$\pi_a=\partial\mathcal L/\partial\dot\phi_a$ 为正则共轭动量，$\mathcal H$ 为哈密顿量密度；空间梯度项取决于 $\mathcal H$ 对 $\nabla\pi_a$、$\nabla\phi_a$ 的依赖。可参考教材式（2.10）的积分形式。
\end{problem}
\begin{answer}
  \begin{equation}
    \frac{\partial \mathcal{L} }{\partial \phi_a } - \partial_\mu \frac{\partial \mathcal{L} }{\partial\left(\partial_\mu\phi_a \right) }
    =0,\quad
    \pi_a
    =\frac{\partial \mathcal{L} }{\partial\dot{\phi}_a },\quad
    \frac{\partial \mathcal{L} }{\partial \phi_a } - \dot{\pi}_a - \partial_i\frac{\partial\mathcal{L} }{\partial\left(\partial_i\phi_a \right) }
    =0. 
  \end{equation}
  一方面，哈密顿量是拉格朗日量的勒让德变换 ：
  \begin{equation}
    \begin{aligned}
      H
      &=\int \mathrm{d}^3\bm{x} \pi_a \dot{\phi}_a - L,
    \end{aligned}
  \end{equation}
  \begin{equation}
    \begin{aligned}
      \delta H
      &=\int \mathrm{d}^3\bm{x} \left(\pi_a \delta\dot{\phi}_a + \dot{\phi}_a\delta \pi_a - \frac{\partial \mathcal{L} }{\partial \phi_a }\delta\phi_a - \frac{\partial \mathcal{L} }{\partial \dot{\phi}_a }\delta\dot{\phi}_a - \frac{\partial \mathcal{L} }{\partial \left(\partial_i \phi_a \right) } \delta\left(\partial_i \phi_a \right)  \right) \\
      &=\int \mathrm{d}^3\bm{x} \left(\dot{\phi}_a\delta \pi_a - \frac{\partial \mathcal{L} }{\partial \phi_a }\delta\phi_a - \frac{\partial \mathcal{L} }{\partial \left(\partial_i \phi_a \right) } \partial_i\left( \delta\phi_a \right)  \right) \\
      &=\int \mathrm{d}^3\bm{x} \left(\dot{\phi}_a\delta \pi_a - \frac{\partial \mathcal{L} }{\partial \phi_a }\delta\phi_a + \partial_i\frac{\partial \mathcal{L} }{\partial \left(\partial_i \phi_a \right) } \left( \delta\phi_a \right)  \right) \\
      &=\int \mathrm{d}^3\bm{x} \left(\dot{\phi}_a\delta \pi_a + \left(\partial_i\frac{\partial \mathcal{L} }{\partial \left(\partial_i \phi_a \right) } - \frac{\partial \mathcal{L} }{\partial \phi_a } \right)\delta\phi_a \right) \\
      &=\int \mathrm{d}^3\bm{x} \left(\dot{\phi}_a\delta \pi_a - \dot{\pi}_a\delta\phi_a \right), \\
    \end{aligned}
  \end{equation}
  其中，分部积分后舍去了边界项。另一方面，$\mathcal{H}=\mathcal{H}(\phi_a,\pi_a,\nabla\phi_a,\nabla\pi_a) $，于是
  \begin{equation}
    \begin{aligned}
      \delta H
      &=\int \mathrm{d}^3\bm{x} \left(\frac{\partial\mathcal{H} }{\partial \phi_a }\delta\phi_a + \frac{\partial\mathcal{H} }{\partial \pi_a }\delta\pi_a + \frac{\partial\mathcal{H} }{\partial (\partial_i \phi_a) }\delta(\partial_i \phi_a) + \frac{\partial\mathcal{H} }{\partial (\partial_i \pi_a) }\delta(\partial_i \pi_a) \right) \\
      &=\int \mathrm{d}^3\bm{x} \left(\frac{\partial\mathcal{H} }{\partial \phi_a }\delta\phi_a + \frac{\partial\mathcal{H} }{\partial \pi_a }\delta\pi_a + \frac{\partial\mathcal{H} }{\partial (\partial_i \phi_a) }\partial_i(\delta \phi_a) + \frac{\partial\mathcal{H} }{\partial (\partial_i \pi_a) }\partial_i(\delta \pi_a) \right) \\
      &=\int \mathrm{d}^3\bm{x} \left(\frac{\partial\mathcal{H} }{\partial \phi_a }\delta\phi_a + \frac{\partial\mathcal{H} }{\partial \pi_a }\delta\pi_a - \partial_i\frac{\partial\mathcal{H} }{\partial (\partial_i \phi_a) }(\delta \phi_a) - \partial_i\frac{\partial\mathcal{H} }{\partial (\partial_i \pi_a) }(\delta \pi_a) \right) \\
      &=\int \mathrm{d}^3\bm{x} \left(\left(\frac{\partial\mathcal{H} }{\partial\phi_a }-\partial_i\frac{\partial\mathcal{H} }{\partial\left(\partial_i\phi_a \right) } \right)\delta\phi_a + \left(\frac{\partial\mathcal{H} }{\partial\pi_a } - \partial_i\frac{\partial\mathcal{H} }{\partial\left(\partial_i\pi_a \right) }   \right)\delta\pi_a   \right),
    \end{aligned}
  \end{equation}
  其中分部积分后也舍去了边界项。二者对比可得：
  \begin{equation}
    \frac{\partial\mathcal{H} }{\partial\pi_a } - \partial_i\frac{\partial\mathcal{H} }{\partial\left(\partial_i\pi_a \right) }
    =\dot{\phi}_a,\quad
    \frac{\partial\mathcal{H} }{\partial\phi_a }-\partial_i\frac{\partial\mathcal{H} }{\partial\left(\partial_i\phi_a \right) }
    =-\dot{\pi}_a,
  \end{equation}
  也即
  \begin{equation}
    \dot{\phi}_a
    =\frac{\partial\mathcal{H} }{\partial\pi_a } - \nabla\cdot\frac{\partial\mathcal{H} }{\partial\left(\nabla\pi_a \right) },\quad
    \dot{\pi}_a
    =-\frac{\partial\mathcal{H} }{\partial\pi_a } + \nabla\cdot\frac{\partial\mathcal{H} }{\partial\left(\nabla\phi_a \right) }. 
  \end{equation}
\end{answer}
\end{document}
